\sscp{Prothesis: irregular C-onsets on V-initial roots}{13-02-02}
\citet[114--118]{Stresemann1927} noticed a number of unexpected initial
consonants on some words in some languages of central Maluku,
and ascribed their development through a process of ‘spiritus
lenis’, particularly in words that begin with {\#}a in nearby languages.
He also observed a similar phenomenon in \ili{Kei},
\tsc{\ili{Teor}-Kur} (\tsc{Seram-Tanimbar-Bomberai} languages), and \ili{Gorontalo} (northern Sulawesi).

\citet{Dyen1962} added an additional sound *W to the \tsc{PAn}
consonant inventory largely in response to this word-initial
phenomenon in data from languages in Maluku and \ili{Chamorro}.
\citet{ChlenovSirk1973} saw an unexpected initial \it{h} in languages
around \ili{Ambon} and Seram and concluded it was a prothetic sound that precedes
the initial vowel (reported in \citealp[121]{Collins1983}).

\citet[6, 120--128, 130f]{Collins1982d,Collins1983}
discusses the same phenomenon but comes to a very different explanation.
He says “various word-initial irregularities in \ili{Central} Maluku
languages can be explained if we assume that these words retained
reflexes of the PAN articles *si and *u” \citep[6]{Collins1983}.
He later concludes “It has been argued that the retention of
\tsc{PAn} *si in some descendants of PCM has resulted in the
appearance of the prothetic segments \it{y-} or \it{h-}.
Possibly a parallel retention of PAN *su or *u is reflected in
the occurrence of initial \it{w-} in certain \ili{Central} Maluku languages” \citep[130]{Collins1983}.
But “the occurrence of prothetic segments in \ili{Central} Maluku languages
cannot constitute evidence in any subgrouping argument involving
\ili{Central} Maluku and Oceania because it is a retention
and not an innovation” \citep[131]{Collins1983}.

Yet Collins also sees subgrouping implications for his explanation:
“In summary two significant phenomena came to light in the course
of this research: the retention of a reflex of PAN *ɫ- distinct
from *l- and the retention of a PAN noun-marker *si
and possibly *su as well.
These two retentions have important ramifications for the subgrouping of PAN.
[{\ldots}] The critical significance of syntactic markers in
PAN as well as in the evolution of its descendants is well attested
in PCM.” (\citet[131]{Collins1983}; PCM = Proto-\ili{Central} Maluku,
see chapters \ref{ch:CM}--\ref{ch:Ban}).
In the data available to us there is even greater variety
and wider distribution than that described above.
We initially focus on \tsc{\ili{Ambon}-Seram} languages
and then expand from there.

It is the restriction of these irregular C-onsets to nouns that
led Collins to posit the retention of \tsc{PAn} nominal marker
*si “personal noun marker” \citep[124]{Collins1983},
which has reflexes found widely in the Philippines as a “personal
subject determiner”, and in western Indonesian languages.
For example, we note that in modern \ili{Malay} \it{si} is commonly
used in the lexicalised interrogative pronoun \it{siapa} ‘who’
(compare \it{apa} ‘what’) and in phrases such as \it{si anu}
or \it{si Polan} ‘what’s-his-name’ (for both terms).
It can be used almost as a kind of definite article before proper
names or nicknames (e.g.
\it{si botak} ‘baldy, the bald one’; or the popular long-running
Indonesian children’s television show \it{Si Unyil}, cf.
\it{unyil} ‘small and cute’), and occasionally to indicate a
kind of personification with other nouns.
Collins describes *su for Proto-Philippines as “the common nominative
determiner” and says it should be assumed that it “can also be reconstructed
in PAN or some language ancestral to both \ili{Central} Maluku
and the Philippine languages” \citep[126]{Collins1983}.
His descriptions vary between *su
and *u.\footnote{
		\citet{Blust2015} reconstructs *su ‘common noun
		case marker’ to Proto-Philippines,
		but is unable to reconstruct it to the level of PMP or \tsc{PAn}.
		\citet{Lichtenberk1988} shows that the kind of prothesis discussed
		in this section is found in some Oceanic languages (thus not
		limited to Wallacea), and may also relate to items with initial PMP *q.}

Collins points out that perhaps not all irregular initial consonants
have developed from this “prothesis” phenomenon.
For example he suggests that the initial {\#}h in some words
in \tsc{East \ili{Piru} Bay} languages are a retention of PMP *q.
The data for \ili{Kamarian},
\ili{Rumakai}, \ili{Haruku}, and \ili{Saparua} in \trf{13-09} are repackaged
and updated from \citet[121]{Collins1983},
and supplemented with our own data from other languages.
The data for ‘nose’ were given by \citet[121]{Collins1983} in
relation to an older reconstruction *[qijuŋ],
whereas \citet{BlustTrussel2020} no longer provide a doublet
for ‘nose’ with a *q onset.
The form for ‘earthquake’ is not found in \citet{BlustTrussel2020}.
These are just a few of the complications that we find in looking
at this issue of irregular consonant onsets.
(Note: words in ‹angle brackets› are from colonial era Dutch
sources that could not be verified by Collins.
In the following three Tables,
light blue shading indicates words with y-onsets; yellow shading
for words with w-onsets; orange for words with \it{l}-onsets;
red for \it{h}-onsets; and grey for other onsets,
such as \it{s-, t-} or \it{v-}.)

\begin{table}\tcp{Possible retention of PMP initial *q in Maluku}{13-09}
\begin{adjustbox}{width=\textwidth}
	\begin{threeparttable}\st{2pt}
	\begin{tabular}{lllllll}\lsptoprule
local gl.		&	pestle									&	liver									&						&	earthquake						&	worm							&	nose\\
PMP					&	*qahəlu									&	*qatay								&						&	*qi(n)sud							&	*quləj						&	*(ŋ/q)ijuŋ\\	\midrule
\ili{Buru}				&	{\hr}alu								&	{\hr}ata-n 						&	‘spleen’	&	{\hr}iso							&	{\hr}ule ‘maggot’	&	\cg{(ŋee-n)}\\
\ili{Kamarian}		&	{\hr}haru					{\crd}&	\cg{(hatuʔa)}					&						&	‹hisu›					{\crd}&	‹hure›			{\crd}&	{\hr}hiru{\crd}\\
\ili{Rumakai}			&	{\hr}haru					{\crd}&	{\hr}ata							&						&	{\hr}ha-hisur-e	{\crd}&	‹hure›			{\crd}&	{\hr}hiru{\crd}\\
\ili{Haruku}			&	‹haru›						{\crd}&	‹hata›					{\crd}&						&	{\hr}isur							&	\cg{(lade)}				&	{\hr}iru\\
\ili{Saparua}			&	‹halu›						{\crd}&	{\hr}ata							&						&	{\hr}isur-o						&	{\hr}ure-ure-o		&	{\hr}irim\\
\ili{Alune}				&													&	{\hr}ata							&						&	{\hr}isu							&	{\hr}ule ‘maggot’	&	{\hr}inu bata-i\\
\ili{Wemale}			&	{\hr}halu					{\crd}&	{\hr}yata				{\cbl}&						&	{\hr}hisu				{\crd}&	{\hr}ulu-te				&	{\hr}ili\\
\ili{Sepa}				&	{\hr}ana alu						&	\cg{(hule)}						&						&	{\hr}isu							&	\cg{(hiayalo)}		&	{\hr}iri-n\\
\ili{Tamilou}			&	\cg{(hatu anan)}				&	\cg{(hule)}						&						&	{\hr}isu							&	\cg{(hiyalo)}			&	{\hr}uli \it{(met.)}\\
\ili{Kowiai}			&	\cg{(lesun-e ana-i)}		&	{\hr}lata\su{†}	{\cor}&						&												&										&	\\
\ili{Yamdena}			&	{\hr} alu								&	{\hr}ata-ŋw						&						&												&	\cg{(kbwatar)}		&	{\hr}iri-ŋw\\
\ili{Kei} Kecil		&	\cg{(luhun yana-n)}			&	{\hr}yata-n			{\cbl}&	‘heart’		&	\cg{(ror)}						&	\cg{(laŋar)}			&	{\hr}niru-n\\
\ili{Fordata}			&	\cg{(saʔa)}							&	{\hr}yata-n			{\cbl}&						&												&	{\hr}ular ‘caterpillar’	&	{\hr}niru-n\\
\ili{Onin}				&													&	{\hr}yata-n			{\cbl}&						&												&		&	\\
\ili{Uruangnirin}	&													&	{\hr}lata-n\su{†}	{\cor}&					&												&		&	\\
		\lspbottomrule
	\end{tabular}
		\begin{tablenotes}\tnsize
			\item[†] {The \it{l}-onsets in \ili{Kowiai}
							and \ili{Uruangnirin} are due to regular **y > \it{l} (see \srf{11-03-02},
							\srf{11-05-01}).}
		\end{tablenotes}
	\end{threeparttable}
\end{adjustbox}
\end{table}

Yet there are also many *q-initial words in these same languages
that show no h-onset,
including reflexes of *qulu ‘head’ (as \it{uru/ulu}),
and *quzan ‘rain’ (as \it{ulan}).
And, as with *ijuŋ in \trf{13-09} (see also \trf{07-12}),
there are words reconstructed without *q which do show h-onsets.
So even if {\#}h is accepted as a retention of PMP *q,
it is lexically specific
and not a regular sound correspondence.
Data from additional languages show that other irregular onsets
occur on these same words regardless of whether or not there
is a retention of PMP *q.

Similarly, \citet{Blust1981} in “The \ili{Soboyo} reflexes of PAN *S”
and \citet{Collins1981} assert that some h-onsets in some words
in \ili{Soboyo} and some languages of \ili{Ambon} are a retention of \tsc{PAn} *S (PMP *h).
Again, this would have to be lexically specific,
since these same languages have other *S-initial words that have no consonant onset.
But being part of the broader pattern of irregular C-onsets found
more widely in Maluku,
as discussed in this section,
provides an equally likely alternate explanation to the rare
and lexically specific retention of \tsc{PAn} *S.

Before discussing Collins’ explanation of irregular consonant
onsets, we provide additional examples in \trf{13-10}
and \trf{13-11} organised by the onset of the historical form.
These Tables present a fuller picture of the data,
including examples from languages that do not show such “prothesis”,
and from a wider distribution of languages that show other onsets
besides y- and w- as well.
Keep in mind that these data are selected to show irregular onsets,
and there are many other words in these languages with the same
historical patterns that do not have these irregular onsets.

\begin{table}[p]\centering\tcp{Irregular C-onsets in Maluku and Bomberai (part 1)}{13-10}
\begin{adjustbox}{width=\textwidth}
	\begin{threeparttable}\st{2pt}
	\begin{tabular}{lll@{ }lllll}\lsptoprule
local gl.		&	dog								&	\mc{2}{l}{thatch}						&	dust, ash								&	calcium						&	jaw, chin					&	wind						\\
PMP					&	*asu							&	\mc{2}{l}{*qatəp}						&	*qabu\su{†}							&	*qapuR						&	*qazay						&	*haŋin					\\	\midrule
\ili{Buru} 				&	{\hr}asu					&	\mc{2}{l}{{\hr}ate-t}				&	{\hr}lafu-n				{\cor}&	{\hr}apu					&	{\hr}aa-n					&	{\hr}aŋin				\\	\midrule
Boano				&										&														&	&	{\hr}lahuʔi				{\cor}&										&	{\hr}ala-					&									\\
\ili{Asilulu}			&	{\hr}asu					&	\mc{2}{l}{\cg{(kaha)}}			&	{\hr}sahut							&	\cg{(kapul)}			&	{\hr}ala-n				&	{\hr}anin				\\
\ili{Haruku}			&	{\hr}asu					&	\mc{2}{l}{{\hr}ate-o}				&	\cg{(hau mata-e)}				&	\cg{(aheti-e)}		&	{\hr}ala					&	{\hr}anin-o			\\
\ili{Saparua}			&	{\hr}asu-ro				&	\mc{2}{l}{{\hr}ato-o}				&													&	{\hr}haullo	{\crd}&	{\hr}wala		{\cye}&	{\hr}anin-o			\\
\ili{Amahai}			&	{\hr}asu-lo				&														&	&	\cg{(hau maa-no)}				&	\cg{(lokal-o)}		&	{\hr}wala		{\cye}&	{\hr}anin-o			\\
\ili{Paulohi}			&	{\hr}asu					&	\mc{2}{l}{{\hr}ate}					&	\cg{(afu mata-i)}				&	\cg{(losa)}				&	{\hr}ala					&	{\hr}anin-e			\\
\ili{Alune}				&	{\hr}asu					&	\mc{2}{l}{{\hr}ate}					&	\cg{(ndopo-ne)}					&	\cg{(asai)}				&	{\hr}ala-mu				&	\cg{(balat-e)}	\\
N. \ili{Wemale}		&	{\hr}yasu		{\cbl}&	{\hr}yate-ye	{\cbl}			&	&	\cg{(mutu-i)}						&	\cg{(kaun-e)}			&	{\hr}yala		{\cbl}&	{\hr}yanin-e	{\cbl}\\
S. \ili{Wemale}		&	{\hr}yasu		{\cbl}&	{\hr}yate			{\cbl}			&	&													&										&	{\hr}ala-mu				&											\\
\ili{Sepa}				&	{\hr}wasu		{\cye}&	\mc{2}{l}{\cg{(laun-o)}}		&	\cg{(hau mata-n)}				&	\cg{(losa)}				&	{\hr}wala-n	{\cye}&	{\hr}yanin-o	{\cbl}\\
\ili{Tehoru}			&	{\hr}wasu		{\cye}&														&	&													&	{\hr}yafi-o	{\cbl}&	{\hr}wala-n	{\cye}&	{\hr}yann-o		{\cbl}\\
\ili{Tamilou}			&	{\hr}wasu		{\cye}&	\mc{2}{l}{\cg{(laun-o)}}		&	\cg{(hau mata-n)}				&	{\hr}haul-o	{\crd}&	{\hr}wala-n	{\cye}&	{\hr}yaann-o	{\cbl}\\
\ili{Yalahatan}		&	{\hr}asu					&	\mc{2}{l}{{\hr}ateaa}				&	{\hr}laho-ne						&	{\hr}ahotoi				&	{\hr}ara-i				&	\cg{(yolil-e)}			\\
S. \ili{Nuaulu}		&	{\hr}asu					&	\mc{2}{l}{{\hr}ainata-e}		&	{\hr}tapu-ne						&										&	{\hr}ana nohu-e		&	\cg{(ihut-e)}				\\
\ili{Manusela}		&	{\hr}wasu-a	{\cye}&	{\hr}wate-a	{\cye}				&	&	{\hr}lahu-a				{\cor}&										&	{\hr}aha-nia			&	{\hr}anino-a				\\
\ili{Liana-Seti}	&	\cg{(inauwa)}			&	\mc{2}{l}{{\hr}ata-lanːe}		&	{\hr}laf-na				{\cor}&	\cg{(loosa)}			&	{\hr}ala-m				&	\cg{(simala)}				\\	\midrule
\ili{Bobot}				&	{\hr}yasw		{\cbl}&	\mc{2}{l}{{\hr}yata}				&													&										&	{\hr}yala-n	{\cbl}&	{\hr}yakin		{\cbl}\\
\ili{Masiwang}		&	{\hr}yai\su{‡}{\cbl}&	\mc{2}{l}{\cg{(ilan)}}		&	{\hr}yafh tai-n					&	{\hr}yavis	{\cbl}&	{\hr}yala-n	{\cbl}&	{\hr}yakin		{\cbl}\\
\ili{Geser}				&	{\hr}asu					&	\mc{2}{l}{\cg{(barém)}}			&	{\hr}afi korus					&	\cg{(lodar)}			&	{\hr}ar						&	{\hr}aŋin						\\
\ili{Kowiai}			&	\cg{(aɸuna)}			&	\mc{2}{l}{\cg{(barem)}}			&	{\hr}abuba							&	{\hr}laɸor	{\cor}&	{\hr}laraʔai{\cor}&	\cg{(lor)}					\\
\ili{Watubela}		&										&	\mc{2}{l}{{\hr}kataf}				&													&										&	\cg{(kakelan)}		&	{\hr}aŋin						\\
\ili{Yamdena}			&	{\hr}asu					&	\mc{2}{l}{{\hr}dafat \it{(met.)}}		&	{\hr}afu, afu-ny				&	{\hr}yafur	{\cbl}&	\cg{(kekan)}			&	\cg{(mnaur)}				\\
\ili{Kei} Kecil		&	{\hr}yahaw	{\cbl}&	\mc{2}{l}{{\hr}rafat \it{(met.)}}		&	{\hr}haruv \it{(met.)}	&	{\hr}yafur	{\cbl}&	\cg{(kakeə-n)}		&	\cg{(niit, yaran)}	\\
\ili{Fordata}			&	{\hr}yaha		{\cbl}&	\mc{2}{l}{{\hr}rafat \it{(met.)}}		&	{\hr}kiavu							&	{\hr}yafur	{\cbl}&	\cg{(demi-n)}			&	\cg{(nait)}					\\
\ili{Sekar}				&	{\hr}yasi		{\cbl}&							&								&													&	{\hr}yafer	{\cbl}&										&	\cg{(dirig)}\\
\ili{Uruangnirin}	&	{\hr}lasi		{\cor}&							&								&													&	{\hr}lafur	{\cor}&										&	\cg{(dirin)}\\
		\lspbottomrule
	\end{tabular}
		\begin{tablenotes}\tnsize
			\item[†]	{\citet{BlustTrussel2020} list *labu ‘ashen,
								grayish’ as a doublet of *qabu ‘ash,
								hearth, cinder, powder, dust; gray’ based on only two witnesses:
								\ili{Buru} \it{lafu} ‘ashes’ and \ili{Kelabit} \it{labuh} ‘grey,
								colour of darker clouds’.
								The \ili{Buru} form only occurs with suffixes \it{lafu-n} and \it{lafu-t}.
								\trf{13-09}--\trf{13-11} show that
								there are several other languages with {\#}l-onsets in this
								and other words, but perhaps not due to the retention of a historical form.}
			\item[‡]	{\ili{Masiwang}: unexplained loss of *s; final *u{\#} > \it{i} is regular.}
		\end{tablenotes}
	\end{threeparttable}
\end{adjustbox}
\end{table}

\trf{13-11} provides additional examples of irregular C-onsets.
There are also further examples of these patterns from additional
languages and etyma in \citet[120--128]{Collins1982d,Collins1983}.

\begin{table}[p]\centering\tcp{Irregular C-onsets in Maluku and Bomberai (part 2)}{13-11}
\begin{adjustbox}{totalheight=\textheight}
	\st{2pt}\begin{tabular}{lllllll}\lsptoprule
*gloss			&	fish gills					&	fire									&	\tsc{1pl.excl}	&	tree, wood			&	rafter							&	fish hook			\\
PMP					&	*hasaŋ							&	*hapuy								&	*kami						&	*kahiw					&	*kasaw							&	**ka[b/w]il		\\	\midrule
\ili{Buru}				&	{\hr}asa-n					&	\cg{(bana)}						&	{\hr}kami				&	{\hr}kau				&	{\hr}kasa						&	{\hr}kawil		\\	\midrule
\ili{Kelang}			&											&	{\hr}au								&									&	{\hr}hai	{\crd}&	{\hr}hasa-e		{\crd}&								\\
\ili{Asilulu}			&	{\hr}asa-na					&	{\hr}au								&	{\hr}ami				&	{\hr}ai					&	{\hr}ʔasa						&	{\hr}ahil			\\
\ili{Alang}				&	{\hr}hasa-na	{\crd}&												&									&									&											&								\\
\ili{Kairatu}			&	{\hr}hasa-i		{\crd}&												&									&									&											&								\\
\ili{Haruku}			&											&	{\hr}hau				{\crd}&	{\hr}ami				&	{\hr}ai					&											&								\\
\ili{Saparua}			&											&	{\hr}hau				{\crd}&	{\hr}ami				&	{\hr}ai					&											&	{\hr}hahillo	{\crd}\\
\ili{Amahai}			&	\cg{(pa-kae)}				&	{\hr}hau				{\crd}&	{\hr}ami				&	\cg{(uwon)}			&	{\hr}aka						&	{\hr}ahir-olo	\\
\ili{Paulohi}			&											&	{\hr}afu							&	{\hr}ami				&	{\hr}ai					&	{\hr}asa						&	{\hr}ahir-e		\\
\ili{Alune}				&	{\hr}sasa-e		{\cdr}&	{\hr}au-e							&	{\hr}ami				&	{\hr}ai, ai-ni	&	{\hr}asakwe					&	{\hr}ahil-e mata-i\\
\ili{Watui}				&											&												&									&									&	{\hr}wasaʔwe	{\cye}&								\\
Naka{\Q}ela	&											&												&									&									&	{\hr}kwasake	{\cye}&								\\
N. \ili{Wemale}		&	{\hr}sasa-e		{\cdr}&	{\hr}yahu				{\cbl}&	{\hr}yami	{\cbl}&	{\hr}yai	{\cbl}&	{\hr}yasa			{\cbl}&								\\
S. \ili{Wemale}		&											&	{\hr}yahu				{\cbl}&									&	{\hr}yai	{\cbl}&	{\hr}yasa-i pul{\cbl}&{\hr}yapil-e	{\cbl}\\
\ili{Sepa}				&											&	{\hr}hau				{\crd}&	{\hr}yam	{\cbl}&	{\hr}yai	{\cbl}&	{\hr}wasa			{\cye}&	{\hr}yahil-o	{\cbl}\\
\ili{Tehoru}			&											&	{\hr}yafu				{\cbl}&	{\hr}yam	{\cbl}&	{\hr}yai	{\cbl}&	{\hr}wasa			{\cye}&	{\hr}yahil-o	{\cbl}\\
\ili{Tamilou}			&											&	{\hr}hau				{\crd}&	{\hr}yam	{\cbl}&	{\hr}yai	{\cbl}&	{\hr}wasa			{\cye}&								\\
\ili{Yalahatan}		&	{\hr}asa-i					&	\cg{(usa)}						&	{\hr}(y)am{\cbl}&	{\hr}ai					&	{\hr}wasa			{\cye}&	{\hr}ahir-e		\\
S. \ili{Nuaulu}		&	{\hr}asa-e					&	\cg{(usa)}						&	{\hr}ami				&	{\hr}ai					&	{\hr}wasa			{\cye}&	{\hr}yahil-o	{\cbl}\\
\ili{Manusela}		&	\cg{(nake-ni)}			&	{\hr}wouw-a			{\cye}&	{\hr}ami				&	{\hr}ai					&	{\hr}hasa			{\crd}&	\cg{(kasaha)}	\\
\ili{Liana-Seti}	&	{\hr}yasa-na	{\cbl}&	{\hr}waha				{\cye}&	{\hr}ʔam				&	{\hr}ai-a				&	{\hr}asa-ʔa					&	{\hr}ail-a		\\	\midrule
\ili{Bobot}				&	{\hr}yasa-n		{\cbl}&	{\hr}yafw				{\cbl}&	{\hr}am					&	{\hr}ai					&	{\hr}asa-						&								\\
\ili{Masiwang}		&	{\hr}yasa-n		{\cbl}&	{\hr}yah, yaaf	{\cbl}&	{\hr}im-				&	{\hr}ai					&	{\hr}asa						&								\\
\ili{Geser}				&											&	{\hr}ahi, afi					&	{\hr}ami				&	{\hr}(k)ai			&											&								\\
\ili{Kowiai}			&	\cg{(laŋgus)}				&	{\hr}laɸ				{\cor}&	{\hr}am					&	{\hr}ai					&	{\hr}asa						&	\cg{(ises)}		\\
\ili{Watubela}		&											&	{\hr}afi							&									&									&											&								\\
\ili{Yamdena}			&	\cg{(ŋeŋ-ar)}				&	{\hr}au								&	{\hr}kam(i)			&	{\hr}ka					&	{\hr}as-e						&	\cg{(irit)}		\\
\ili{Kei} Kecil		&	{\hr}vaha-n		{\cbl}&	{\hr}yaf				{\cbl}&	{\hr}ʔam				&	{\hr}ay					&	{\hr}aha						&	{\hr}ail			\\
\ili{Fordata}			&	{\hr}ala-n					&	{\hr}yafu				{\cbl}&	{\hr}ami				&	{\hr}aa					&	{\hr}aha						&	\cg{(ihi)}		\\
\ili{Onin}				&											&	{\hr}yafi				{\cbl}&	{\hr}ami				&									&											&								\\
\ili{Sekar}				&											&	{\hr}yafi				{\cbl}&	{\hr}yami	{\cbl}&									&											&								\\
\ili{Uruangnirin}	&											&	{\hr}lafi				{\cor}&	{\hr}ami				&									&											&								\\
\ili{Arguni}			&											&	{\hr}iyáf				{\cbl}&	{\hr}mbami{\cdr}&									&											&								\\
\ili{Bedoanas}		&											&	{\hr}iaf				{\cbl}&	{\hr}mbami{\cdr}&									&											&								\\
\ili{Erokwanas}		&											&	{\hr}iaf				{\cbl}&	{\hr}mbemi{\cdr}&									&											&								\\
		\lspbottomrule
	\end{tabular}
\end{adjustbox}
\end{table}

The onset of the form \it{hau} ‘fire’ in several languages from
*hapuy ‘fire’ warrants discussion.
The \it{h}-onset could be: (a) a retention of PMP *h; (b) metathesis
involving medial *p; or (c) one of these irregular onsets that
we are talking about in this section.
\citet[115, 118]{Collins1983} argues for metathesis of *p.
We note that: 1) the \it{hau} ‘fire’ form is found in languages
where regularly *p > \it{h},
so metathesis is possible;\footnote{
		Data from 13 nearby languages
		not presented here have the form \it{au},
		but in those languages *p > {\0},
		so the data from these languages is ambiguous as to whether or
		not metathesis has taken place.}
2) different onsets
(\it{{\#}y} and \it{{\#}l}) are found in eight other languages for this same
etymon, indicating this is one of those words that triggers or allows such onsets.
So while we acknowledge that metathesis is a possible explanation,
regularity of form, function
and distribution mean that \it{{\#}h} cannot be ruled out here
as an irregular onset.

We also note that the similarity in form between *hapuy ‘fire’
and *qabu ‘ash’, along with related meaning,
mean that distinguishing derivative forms resulting from each
of these etyma is not always clear.
Several patterns can be seen from the data in \trf{13-09}--\trf{13-11}:

\begin{itemize}
		\item	Some languages in Maluku
					and Bomberai have irregular onsets with \it{y}-insertion,
					\it{w}-insertion, \it{h}-insertion, \it{l}-insertion, or \it{s}-insertion
					(in descending order of frequency),
					along with a few other unexplained onsets.
		\item	We notice a slight correlation between words that have
					\it{w}-onsets with historical forms that have *u or *w (or *-aw)
					within them, but also see enough exceptions to only observe a tendency,
					rather than posit a conditioning rule.
		\item	These irregular onsets are mostly before nouns
					and pronouns that begin with {\#}a in nearby languages.
					Other data not presented here have possible candidates for irregular
					onsets with words from other parts of speech
					and with different vowels (e.g.
					PMP *hasaq ‘whet, sharpen’ (verb),
					*pija ‘how much, how many’ (question word,
					quantifier), *pəñu ‘sea turtle’).
		\item	Individual etyma (the vertical columns) can show different
					onsets in different languages.
					So *asu ‘dog’ has \it{{\#}a, {\#}y, {\#}w};
					*qabu, ‘dust,
					ash’ has \it{{\#}a, {\#}l, {\#}s, {\#}t, {\#}y};
					*qapuR ‘calcium, chalk,
					mineral lime’ has \it{{\#}a, {\#}h, {\#}y, {\#}l};
					*qazay ‘jaw, chin’ has
					\it{{\#}a, {\#}w, {\#}y, {\#}l}, and so forth.
		\item	Some individual languages (the horizontal rows) seem to
					use a single onset for those words that take one.
					For example, in our data \ili{Wemale},
					\ili{Bobot}, and \ili{Masiwang} seem to use \it{{\#}y} on (almost) all words
					that have irregular onsets.
					And \ili{Kowiai} and \ili{Uruangnirin} use \it{{\#}l} on all words that have
					irregular onsets, but this is due to a regular **y > \it{l}
					sound change in \ili{Uruangnirin} (cf.
					*buqa\tbr{y}a > \it{pua\tbr{l}a} ‘crocodile’).
		\item	Other individual languages,
					such as \ili{Sepa}, \ili{Tehoru},
					and \ili{Tamilou} (\tsc{East Littoral},
					\srf{10-01-02}), seem to enjoy a variety of different onsets for different words.
					Although some clusters of languages seem to work together as
					a packet, such as North and South \ili{Wemale}, or \tsc{East Littoral},
					we also see internal variation within those clusters.
\end{itemize}

We note that most of the languages that have irregular onsets
on nouns also have subject indexing on active verbs,
or various stative prefixes on non-active verbs,
such that bare verb roots are rarely found as onsets of words.
This may provide a partial explanation why such irregular onsets
do not appear to be found on verbs,
and seem to be restricted to nouns.

We have several concerns regarding Collins’ explanation of irregular
consonant onsets that “various word-initial irregularities in
\ili{Central} Maluku languages can be explained if we assume that these
words retained reflexes of the PAN articles *si and *u” \citep[6]{Collins1983}.

First of all, the functions are different.
Where the *si marker occurs unambiguously in western Austronesian
languages with form and function intact, it always involves the identity of people or signals personification.
None of the data under scrutiny that we have access to in this
region involves the identity of people.
The only thing in common is that irregular onsets happen to nouns -- albeit
very different kinds of nouns not involving personal identity.
\citet[124]{Collins1983} feels he can get around this by declaring,
“In these [\tsc{East \ili{Piru} Bay}] languages,
*si was not restricted to personal names or pronouns; it occurred
with all nouns.” 

Secondly, to accept Collins’ argument,
the forms found in languages of Maluku would require us to abandon
common principles of regular sound correspondences.
These languages regularly retain reflexes of PMP *s (see reflexes
of *qisud in \trf{13-09}; *asu in \trf{13-10}; *hasaŋ
and *kasaw in \trf{13-11}).
Other data not presented here show that reflexes of PMP *s are
also retained in the initial position, and before *i.
Etyma we checked for these languages were *sakay ‘ascend’,
*saŋa ‘bifurcation, fork’, *səpsəp ‘suck’, *suqan ‘digging stick’,
*suja ‘spike trap’, *susu ‘female br\ili{east}’, *siwa ‘nine’,
*ŋi(n)si ‘grin, show teeth (local gl.: teeth)’.

Collins is aware of this problem
and tries to get around it by declaring “we note that PAN *s
is regularly retained in all E[ast] P[iru] B[ay] languages.
Here we claim that *si- became \it{h-} or \it{y-}.
The fact that *si- was a morphological marker made it possible
for *si- to undergo innovations which did not affect *si in other reflexes.
Special treatment of morphological affixes is
well known in other language families.” \citep[124f]{Collins1983}.

We find Collins’ proposal of retentions of *si
and *(s)u inadequate to explain:

\begin{itemize}
	\item	the great variety of onsets found collectively in the languages,
				with y-insertion, w-insertion, h-insertion, l-insertion, s-insertion
				along with some other possibilities such as t-insertion and v-insertion.
	\item	the variety of onsets found in single languages.
	\item	why, for the same etymon,
				one language would have no onset,
				a nearby related language would have y-insertion,
				another nearby language w-insertion,
				and another h-insertion or l-insertion for the same word.
\end{itemize}

We find Collins’ proposal unconvincing that \it{w}-onsets in
\tsc{East \ili{Piru} Bay} languages in Maluku reflect a *su nominal marker
reconstructed for Proto-Philippines, and that therefore *su should be reconstructed at a level ancestral
to both Proto-Philippines and \ili{Central} Maluku.
It only attempts to address a small portion of the data
and fails to explain the whole phenomenon of many irregular onsets.
We notice that Collins himself is often a bit hesitant about *su \citep[130--131]{Collins1983}.

In \trf{13-09}, \trf{13-10}, and \trf{13-11}, these irregular
onsets are found sporadically in several subgroups: \tsc{Sula-Buru}, \tsc{\ili{Ambon}-Seram}, and \tsc{Seram-Tanimbar-Bomberai}.
In \tsc{Timor-Babar} \ili{Idate}, irregular \it{y}-onsets are found -- some
involving the same words that have irregular onsets in \tsc{\ili{Ambon}-Seram}
(see Tables \ref{tab:04-01}, \ref{tab:04-09}, \ref{tab:04-11}).
Tables \ref{tab:04-33} and \ref{tab:04-34} show irregular \it{w}-onsets for \tsc{Timor-Babar}
languages \tsc{\ili{Teun}-\ili{Nila}-Serua} -- some involving the same words
that have irregular onsets in \tsc{\ili{Ambon}-Seram}.
\trf{11-27} presents additional examples of \it{y}-onsets in
\tsc{Seram-Tanimbar-Bomberai} languages for *anak ‘child’,
*ama ‘father’, and *halas ‘forest’, beyond what is presented
earlier in this section.
Section \srf{06-01-03} describes irregular initial glide insertion
in \tsc{Aru} languages (usually with subsequent glide fortition),
some involving the same words that have irregular onsets in \tsc{\ili{Ambon}-Seram},
and many of which involve other vowels.

Irregular onsets are occasionally relevant for lower level subgrouping
elsewhere, as discussed in the relevant chapters.
Specifically, the \it{y}-onsets in \tsc{Nuclear Tanimbar-Bomberai}
are systematic enough to be useful for lower-level subgrouping (\srf{11-05-01}, particularly \trf{11-28}).
Likewise, we find the onsets in \tsc{Timor-Babar} languages to be systematic
enough to be useful for subgrouping at lower levels for \tsc{\ili{Teun}-\ili{Nila}-Serua}
(\srf{04-06-04}, particularly \trf{04-40}), and for subgrouping in \tsc{Aru} (\srf{06-01-03}).
We do not find the irregular onsets in \tsc{East \ili{Piru} Bay} languages
in \tsc{\ili{Ambon}-Seram} to be systematic enough to be convincing
for subgrouping purposes, except perhaps at the lowest levels,
and suggest the classification based on elusive *si
and *su \citep[128]{Collins1983} needs to be revisited in light
of the fuller data
and more widespread phenomenon.

In \tsc{Sula-Buru}, \ili{Buru} and \ili{Lisela} are the only languages that
show a \it{y}-onset in the \tsc{1sg} pronoun \it{yako, ya}.
Since the other \tsc{Sula-Buru} languages (\ili{Mangole} \it{aku},
\ili{Sanana} \it{ak}, \ili{Ambelau} \it{au-ni}) do not have such an onset,
the question naturally arises,
is this a retention of PMP *i-aku ‘\tsc{1sg},
I, me’, or is it one of these irregular onsets found widely in
this region? (Compare PMP *kami > \it{yam(i)} ‘\tsc{1pl.excl}’
in many languages in \trf{13-08}.) Pronominal forms with initial
*i- are found throughout the Austronesian world including in
Sulawesi, and are reconstructed to \tsc{PAn} *i-aku ‘\tsc{1sg}’
and *i-ami ‘\tsc{1pl.excl}’ (\citealp{BlustTrussel2020}, \citealp[908]{AdelaarHajek2024}).
In \tsc{Sula-Buru}, \tsc{\ili{Ambon}-Seram}, and \tsc{Seram-Tanimbar-Bomberai}
subgroups, languages with \it{{\#}y/i} in their pronouns are
interspersed among languages without \it{{\#}y/i},
so it is not clear if we are dealing with retentions or this sporadic prothesis phenomenon.
But the presence of other onsets in these same pronouns suggests
the latter is involved in at l\ili{east} some languages.
Regardless of whether one sees these irregular onsets as retentions
(as per Collins), or as sporadic innovations,
they appear to have little value for higher-level subgrouping.
